% =====================================================================
% Abstract -- DRAFT (2026-10-07)
% ---------------------------------------------------------------------
% Sources: MDPI Sensors paper (Experiment 1), chapter4_methodology.tex,
% chapter7_experiment3_reconstruction.tex and the RIoT v5 results report
% (docs/experiments/riot_v5/riot_v5_results_report.md).
% Lines marked "% AUTHOR-CHECK" need confirming before submission.
% All v5 figures are single-seed (seed 42, 103 runs); re-check them once
% the tier-3 seeds (43/44) have run.
% =====================================================================

\begin{abstract}
Low-cost sensor networks give cities a near-real-time picture of how
people move, but the data they produce are fragmented: sensors drop out
for minutes or days, and a missing interval is easily mistaken for an
absence of activity. This thesis investigates how spatio-temporal deep
learning can deliver reliable pedestrian forecasts for operational
situational awareness when its inputs are incomplete, through three
progressive experiments on the Newcastle Urban Observatory
pedestrian-counting network.

Experiment~1 forecasts 110 sensor streams individually and shows that
the structure of missingness, not its volume, governs accuracy. Error
rises sharply once the mean uninterrupted sequence falls below about
three hours, and completeness has no independent effect once sequence
length is controlled for. Experiment~2 adds graph-structured spatial
modelling across 54 stations
% AUTHOR-CHECK: add the one-sentence headline result of Experiment 2
% (Chapter 6); it was not available when this draft was written.
[\textit{key result}]. Experiment~3 develops RIoT, a multi-task
architecture that couples a graph learnt from urban morphology, GATv2
spatial attention and a Transformer temporal encoder with joint
forecasting, imputation and zero-shot kriging heads. It is trained on
three years of 15-minute counts from 81 stations, 24 of which are
withheld entirely to test transfer to unseen locations.

A controlled study of 103 configurations, scored on observed values
under a leakage-free chronological protocol, shows that the original
formulation, which treated missing targets as zeros, learnt to forecast
zero during outages, and that its evaluation metric rewarded this
collapse. Training on observed targets removes the collapse. The
principal resilience mechanism is a unified missing-data encoding that
makes artificial and genuine gaps indistinguishable to the model,
combined with structured artificial gaps during training. The final
model reduces forecast error under synthetic 3--24-hour outages by
50\% and at stations seen in training by 23\%, and keeps predicted
volume during day-long outages at the true level rather than near zero.
% AUTHOR-CHECK: 50% = SO@12 RMSE 31.8 -> 15.8 (V4R -> V5F); 23% =
% transductive RMSE 17.0 -> 13.1; SO@96 mass ratio 0.00 -> 1.01.
Removing the unified encoding alone reverses most of this gain. Graph
message passing matters chiefly during long outages. The results also
mark clear limits: at unseen stations neither the forecaster nor the
kriging head outperforms a neighbour-based seasonal profile, and every
model under-predicts the highest counts. Several earlier conclusions,
including the benefit of auxiliary imputation and the weakness of a
recurrent baseline, prove to be artefacts of how missing data were
scored.

The thesis contributes an empirical account of how data quality limits
urban forecasting, an outage-resilient forecasting architecture whose
robustness is traced to specific ablated components, and a
demonstration that the scoring of missing data can invert the
conclusions of a forecasting study.
\end{abstract}
